\usepackage[margin=1cm,a5paper]{geometry}

\usepackage{bold-extra}
\usepackage{polski}
\usepackage{iwona}
\usepackage[OT4]{fontenc}

\usepackage{amsmath,amssymb,amsthm}
\usepackage{pgf,tikz}
\usepackage{hyperref}

\pagestyle{empty}
\setlength{\parindent}{0pt}
\setlength{\parskip}{5pt}
\renewcommand\leq\leqslant
\renewcommand\geq\geqslant

\colorlet{maincolor}{ForestGreen}
\usepackage[most]{tcolorbox}
\usepackage{varwidth}

\newtcolorbox{tresc}[3][]{enhanced, sharp corners, boxrule=0mm, colback=maincolor!10,
colframe=black,toptitle=1mm,fonttitle=\bfseries,
colbacktitle=maincolor!25, coltitle=black,
title={Zadanie #2\hfill{\it Opracowanie: #3}},#1}

\newenvironment{problem}[3]{\newpage\begin{tresc}{#1}{#2}{#3}}{\end{tresc}}

% Dodatkowe makra

\usepackage{multicol}
\usepackage{cancel}
\usepackage{bbm}

\def\ind{\mathbbm{1}}

\def\N{\mathbb{N}}
\def\P{\mathbb{P}}
\def\E{\mathbb{E}}
\def\Z{\mathbb{Z}}
\def\Q{\mathbb{Q}}
\def\R{\mathbb{R}}

\def\nwd{\mathrm{NWD}}

\newtheorem*{lemat}{Lemat}

\usepackage{enumitem}
\setlist[enumerate,1]{label={(\alph*)},topsep=0pt}

% ------------------------------------------------------------------
% Kolorowe boxy: definicje (zielone), twierdzenia/obserwacje/wnioski (niebieskie)
% ------------------------------------------------------------------

\colorlet{defcolor}{ForestGreen}
\colorlet{thmcolor}{RoyalBlue}

\newcounter{defcnt}
\newcounter{thmcnt}

% wspólny styl boxa
\tcbset{notatkibox/.style={enhanced, sharp corners, boxrule=0mm, frame hidden,
  toptitle=0.5mm, bottomtitle=0.5mm, left=2mm, right=2mm, top=1.5mm, bottom=1.5mm,
  fonttitle=\bfseries, coltitle=black, breakable, before skip=6pt, after skip=6pt}}

% \begin{definicja}[nazwa]
\NewDocumentEnvironment{definicja}{o}{%
  \refstepcounter{defcnt}%
  \begin{tcolorbox}[notatkibox, colback=defcolor!8, colframe=defcolor!25,
    colbacktitle=defcolor!25,
    title={Definicja \thedefcnt\IfValueT{#1}{\ (#1)}}]%
}{\end{tcolorbox}}

% generator środowisk niebieskich: #1 = nazwa środowiska, #2 = etykieta
\newcommand{\newthmbox}[2]{%
  \NewDocumentEnvironment{#1}{o}{%
    \refstepcounter{thmcnt}%
    \begin{tcolorbox}[notatkibox, colback=thmcolor!8, colframe=thmcolor!25,
      colbacktitle=thmcolor!25,
      title={#2 \thethmcnt\IfValueT{##1}{\ (##1)}}]%
  }{\end{tcolorbox}}%
}

\newthmbox{twierdzenie}{Twierdzenie}
\newthmbox{obserwacja}{Obserwacja}
\newthmbox{wniosek}{Wniosek}
\newthmbox{stwierdzenie}{Stwierdzenie}

% treść zadania -- żółty box, \begin{zadanie}{1.1}
\colorlet{zadcolor}{Dandelion}

\NewDocumentEnvironment{zadanie}{m}{%
  \begin{tcolorbox}[notatkibox, colback=zadcolor!15, colframe=zadcolor!40,
    colbacktitle=zadcolor!40,
    title={Zadanie #1}]%
}{\end{tcolorbox}}

% przykład -- bez boxa, wyróżniony nagłówkiem
\theoremstyle{definition}
\newtheorem*{przyklad}{Przykład}

\renewcommand{\qedsymbol}{$\square$}
